\newpage
\section{Manifold Unscented Kalman Filter}

\subsection{Kalman Filter Overview}
The Kalman filter (KF) is an optimal recursive estimator used to estimate a 
state under noisy measurements and uncertain dynamics. It works under an assumption that the relationship
between our state and measurements can be modelled as linear and Gaussian \fixme{cite intro to KF}.

\begin{table}[!htbp]
\centering
\small
\begin{tabular}{p{3.5cm} p{8.5cm}}
\hline
\textbf{Term} & \textbf{Description} \\
\hline
State Vector ($\mathbf{x}$) & The quantity the filter is trying to estimate. \\
State Estimate Covariance ($\mathbf{P}$) & The uncertainty the filter has in its state estimation. \\
Process Noise Covariance ($\mathbf{Q}$) & The uncertainty the filter has in its process model, i.e. how much the filter trusts the dynamics used to propagate the state. \\
Measurement Noise Covariance ($\mathbf{R}$) & The uncertainty the filter has in its measurement values, used to correct the state. \\
Kalman Gain ($\mathbf{K}$) & Balances the weight given to the process model prediction versus the incoming measurement. \\
\hline
\end{tabular}
\caption{Kalman Filter Terms and Definitions}
\label{tab:kf_terms}
\end{table}

The filter implements a two step process of prediction and update. The predicition
propagates our state estimation and the update corrects the state estimation using one
or multiple measurements. The covariance P is updated by the K after every new sensor value to reflect
the new information the filter has gained about the system.

However, the KF breaks down for orientation estimation due to some key factors. First, quaternion kinematics are
non-linear (\fixme{see SO(3) equation}) and the assumption that the relationship is linear is violated. Second, as the relationship is non-linear, this
breaks the Gaussian noise process model as propagation through a non-linear function produces a non-Gaussian output. Third,
as quaternions exist in SO(3) and not Euclidean space, adding a Gaussian term to the state can push it off the manifold.




\subsection{Unscented Kalman Filter}
The Unscented Kalman Filter (UKF) is an extension of the KF that allows for non-linear relationships to be more accurately modelled by the filter. 
Traditionally, to apply a KF to a non-linear relationship an extended KF (EKF) is used which linearises all non-linear 
models through the use of Jacobian matrices that are evaluated at the current state estimate. \fixme{cite Julier}. 
However, better performance tends to be achieved by a UKF for orientation estimation as a UKF gives more accurate covariance propagation at 
larger rotations \cite{kraft_2003_a} \fixme{cite Julier}.

Sigma points are drawn from a specific, deterministic algorithm that when a non-linear function is applied to these points; a set of $2n + 1$ transformed
points are left ($n = 3$ for a 3D error state, giving $7$ sigma points).
These sigma points are then used to reconstruct the mean by taking a weighted average of these points and the weighted outer product
for the covariance. 

However, with the non-linear propagation addressed, computing a weighted mean of quaternions directly is geometrically invalid. This is due to
quaternions exisiting on the SO(3) manifold, as the element-wise average of quaternions does not lie on the unit sphere and therefore does not represent a valid rotation.  

For an in-depth derivation regarding the UKF see \fixme{cite Julier}.

\subsection{The Manifold Problem}
The core issue of the UKF is that it assumes that the x lies within Euclidean space, however quaternions lie on the SO(3) manifold. So computing the mean
and adding the measurement noise to the state is geometrically invalid, thus the state that is estimated is also invalid.

The solution to this problem is to perform any arithmetic computations that the UKF requires within the local tangent space (so(3)). By taking the current 
state estimate (quaternion on the SO(3) manifold) and mapping it to and from Euclidean vector space, the predict and update steps are performed validly 
and the geometric representation of the state is conserved. Brossard et al and Crassidis et al
explore and show this within their UKF-Manifold and Unscented Quaternion Estimator (USQUE) implementations \fixme{cite UKF-M and USQUE}.

Thus the implementation used in this project is based on UKF-M as it is generalisable to various manifolds not just SO(3).

\subsection{Filter Design}



\subsection{Filter Setup}

Table 1 below shows the configuration used for the UKF for each of the scenarios. 

\label{sec:filtersetup}
\begin{table}[!htbp]
\centering
\small
\begin{tabular}{p{2.0cm} p{5.0cm} p{5.0cm} p{1.5cm}}
\hline
\textbf{State} & \textbf{Process Noise ($Q$)} & \textbf{Measurement Noise ($R$)} & \textbf{Sampling Rate} \\
\hline
\makecell[l]{$q_w, q_x, q_y, q_z,$ \\ $\omega_{b,x}, \omega_{b,y}, \omega_{b,z}$}
&
Gaussian noise determined by rest phases
&
Gaussian noise determined by rest phases
&
$286\,\text{Hz}$ \\
\hline
\end{tabular}
\caption{UKF configuration}
\label{tab:filtersetup}
\end{table}

The state estimates the orientation that is determined by integrating gyroscope measurements
where there is an assumed slowly varying additive bias applied to the gyroscope. The matrix $Q$ handles the uncertainty
the filter has in the process model (gyro integration). The measurement model assumes that accelerometer updates
are the measured specific force that aligns with gravity (low linear acceleration), and the magnetometer is a reference
of the local Earth's field. The matrix $R$ is the uncertainty that the filter has in accelerometer and magnetometer
updates and is adaptively weighted.   

\subsubsection{Adaptive Gating for Disturbed Measurements}

The matrix $R$ applied an adaptive scheme where disturbed accelerometer and 
magnetometer measurements were downweighted for correcting the orientation. The initial covariance
was determined by taking a mean of all measurements for each sensor during the rest periods, 
and subtracting this mean, to get residuals from the sensors. These
residuals were then computed as the covariance which is modelled as Gaussian noise and before each correction applied by the filter, 
a heuristic approach was taken by tuning arbitrary tolerances to accept or downweight the sensor
measurements. If measurements were deemed unreliable, then the matrix $R$ would be multiplied by a scalar
value to downweight updates during corrections, this scheme was selected because outright rejection of measurements caused 
discontinuous orientation estimates.

The matrix $Q$, also took a similar approach where gyro bias was determined through subtracting 
the mean during rest phases and these residuals were computed as the covariance which was modelled to be Gaussian noise. 
\newpage


